% Options for packages loaded elsewhere
% Options for packages loaded elsewhere
\PassOptionsToPackage{unicode}{hyperref}
\PassOptionsToPackage{hyphens}{url}
%
\documentclass[
  ignorenonframetext,
  aspectratio=169,
  russian,
]{beamer}
\newif\ifbibliography
\usepackage{pgfpages}
\setbeamertemplate{caption}[numbered]
\setbeamertemplate{caption label separator}{: }
\setbeamercolor{caption name}{fg=normal text.fg}
\beamertemplatenavigationsymbolshorizontal
% remove section numbering
\setbeamertemplate{part page}{
  \centering
  \begin{beamercolorbox}[sep=16pt,center]{part title}
    \usebeamerfont{part title}\insertpart\par
  \end{beamercolorbox}
}
\setbeamertemplate{section page}{
  \centering
  \begin{beamercolorbox}[sep=12pt,center]{section title}
    \usebeamerfont{section title}\insertsection\par
  \end{beamercolorbox}
}
\setbeamertemplate{subsection page}{
  \centering
  \begin{beamercolorbox}[sep=8pt,center]{subsection title}
    \usebeamerfont{subsection title}\insertsubsection\par
  \end{beamercolorbox}
}
% Prevent slide breaks in the middle of a paragraph
\widowpenalties 1 10000
\raggedbottom
\AtBeginPart{
  \frame{\partpage}
}
\AtBeginSection{
  \ifbibliography
  \else
    \frame{\sectionpage}
  \fi
}
\AtBeginSubsection{
  \frame{\subsectionpage}
}
\usepackage{iftex}
\ifPDFTeX
  \usepackage[T1]{fontenc}
  \usepackage[utf8]{inputenc}
  \usepackage{textcomp} % provide euro and other symbols
\else % if luatex or xetex
  \usepackage{unicode-math} % this also loads fontspec
  \defaultfontfeatures{Scale=MatchLowercase}
  \defaultfontfeatures[\rmfamily]{Ligatures=TeX,Scale=1}
\fi
\usepackage{lmodern}

\ifPDFTeX\else
  % xetex/luatex font selection
\fi
% Use upquote if available, for straight quotes in verbatim environments
\IfFileExists{upquote.sty}{\usepackage{upquote}}{}
\IfFileExists{microtype.sty}{% use microtype if available
  \usepackage[]{microtype}
  \UseMicrotypeSet[protrusion]{basicmath} % disable protrusion for tt fonts
}{}


\usepackage{longtable,booktabs,array}
\usepackage{calc} % for calculating minipage widths
\usepackage{caption}
% Make caption package work with longtable
\makeatletter
\def\fnum@table{\tablename~\thetable}
\makeatother
\usepackage{graphicx}
\makeatletter
\newsavebox\pandoc@box
\newcommand*\pandocbounded[1]{% scales image to fit in text height/width
  \sbox\pandoc@box{#1}%
  \Gscale@div\@tempa{\textheight}{\dimexpr\ht\pandoc@box+\dp\pandoc@box\relax}%
  \Gscale@div\@tempb{\linewidth}{\wd\pandoc@box}%
  \ifdim\@tempb\p@<\@tempa\p@\let\@tempa\@tempb\fi% select the smaller of both
  \ifdim\@tempa\p@<\p@\scalebox{\@tempa}{\usebox\pandoc@box}%
  \else\usebox{\pandoc@box}%
  \fi%
}
% Set default figure placement to htbp
\def\fps@figure{htbp}
\makeatother



\ifLuaTeX
\usepackage[bidi=basic,provide=*]{babel}
\else
\usepackage[bidi=default,provide=*]{babel}
\fi
% get rid of language-specific shorthands (see #6817):
\let\LanguageShortHands\languageshorthands
\def\languageshorthands#1{}


\setlength{\emergencystretch}{3em} % prevent overfull lines

\providecommand{\tightlist}{%
  \setlength{\itemsep}{0pt}\setlength{\parskip}{0pt}}



 

\usepackage[]{csquotes}

% _resources/tex/beamer.tex
\usepackage{fontspec}
\usepackage{polyglossia}

\setmainlanguage{russian}
\setotherlanguage{english}

% Use DejaVu fonts (excellent Cyrillic support, usually pre-installed)
\setmainfont{DejaVu Serif}[Ligatures=TeX]
\setsansfont{DejaVu Sans}[Ligatures=TeX]
\setmonofont{DejaVu Sans Mono}

%%% Keep your existing theme selection
\IfFileExists{beamerthemegotham.sty}{
  \usetheme{gotham}
  \gothamset{
    numbering=totalpagenumber,
    parttocframe default=off,
    sectiontocframe default=off,
    subsectiontocframe default=off,
  }
}{
  \usetheme{Madrid}
}
\makeatletter
\@ifpackageloaded{caption}{}{\usepackage{caption}}
\AtBeginDocument{%
\ifdefined\contentsname
  \renewcommand*\contentsname{Содержание}
\else
  \newcommand\contentsname{Содержание}
\fi
\ifdefined\listfigurename
  \renewcommand*\listfigurename{Список иллюстраций}
\else
  \newcommand\listfigurename{Список иллюстраций}
\fi
\ifdefined\listtablename
  \renewcommand*\listtablename{Список таблиц}
\else
  \newcommand\listtablename{Список таблиц}
\fi
\ifdefined\figurename
  \renewcommand*\figurename{Рисунок}
\else
  \newcommand\figurename{Рисунок}
\fi
\ifdefined\tablename
  \renewcommand*\tablename{Таблица}
\else
  \newcommand\tablename{Таблица}
\fi
}
\@ifpackageloaded{float}{}{\usepackage{float}}
\floatstyle{ruled}
\@ifundefined{c@chapter}{\newfloat{codelisting}{h}{lop}}{\newfloat{codelisting}{h}{lop}[chapter]}
\floatname{codelisting}{Список}
\newcommand*\listoflistings{\listof{codelisting}{Листинги}}
\makeatother
\makeatletter
\makeatother
\makeatletter
\@ifpackageloaded{caption}{}{\usepackage{caption}}
\@ifpackageloaded{subcaption}{}{\usepackage{subcaption}}
\makeatother

\usepackage{bookmark}
\IfFileExists{xurl.sty}{\usepackage{xurl}}{} % add URL line breaks if available
\urlstyle{same}
\hypersetup{
  pdftitle={Лабораторная работа №15},
  pdfauthor={бахи сиди али темассини},
  pdflang={ru-RU},
  hidelinks,
  pdfcreator={LaTeX via pandoc}}


\title{Лабораторная работа №15}
\subtitle{Администрирование локальных сетей}
\author{бахи сиди али темассини}
\date{2026-05-23}

\begin{document}
\frame{\titlepage}

\renewcommand*\contentsname{Содержание}
\begin{frame}[allowframebreaks]
  \frametitle{Содержание}
  \setcounter{tocdepth}{1}
  \tableofcontents
\end{frame}
\setcounter{tocdepth}{1}
\tableofcontents
}

\section{Цель
работы}\label{ux446ux435ux43bux44c-ux440ux430ux431ux43eux442ux44b}

\begin{frame}{Цель работы}
\begin{itemize}
\tightlist
\item
  Настройка динамической маршрутизации между территориями организации
\end{itemize}
\end{frame}

\section{Выполнение лабораторной
работы}\label{ux432ux44bux43fux43eux43bux43dux435ux43dux438ux435-ux43bux430ux431ux43eux440ux430ux442ux43eux440ux43dux43eux439-ux440ux430ux431ux43eux442ux44b}

\begin{frame}[fragile]{Настройка OSPF на маршрутизаторе сети «Донская»}
\phantomsection\label{ux43dux430ux441ux442ux440ux43eux439ux43aux430-ospf-ux43dux430-ux43cux430ux440ux448ux440ux443ux442ux438ux437ux430ux442ux43eux440ux435-ux441ux435ux442ux438-ux434ux43eux43dux441ux43aux430ux44f}
\begin{itemize}
\tightlist
\item
  Включён процесс OSPF на маршрутизаторе сети «Донская»
\item
  Назначен Router ID \texttt{10.128.254.1}
\item
  Сети \texttt{10.0.0.0/8} добавлены в \texttt{area\ 0}
\end{itemize}
\end{frame}

\begin{frame}
\begin{figure}

\centering{

\includegraphics[width=0.7\linewidth,height=\textheight,keepaspectratio]{image/1.png}

}

\caption{\label{fig-1}Настройка OSPF на маршрутизаторе сети «Донская»}

\end{figure}%
\end{frame}

\begin{frame}[fragile]{Проверка состояния OSPF и таблицы маршрутизации}
\phantomsection\label{ux43fux440ux43eux432ux435ux440ux43aux430-ux441ux43eux441ux442ux43eux44fux43dux438ux44f-ospf-ux438-ux442ux430ux431ux43bux438ux446ux44b-ux43cux430ux440ux448ux440ux443ux442ux438ux437ux430ux446ux438ux438}
\begin{itemize}
\tightlist
\item
  Выполнена проверка состояния OSPF
\item
  Подтверждён запуск процесса OSPF
\item
  Проверена таблица маршрутизации
\item
  Обнаружен маршрут к сети Сочи \texttt{10.130.0.0/16}
\end{itemize}
\end{frame}

\begin{frame}
\begin{figure}

\centering{

\includegraphics[width=0.7\linewidth,height=\textheight,keepaspectratio]{image/2.png}

}

\caption{\label{fig-2}Проверка состояния OSPF и таблицы маршрутизации}

\end{figure}%
\end{frame}

\begin{frame}[fragile]{Настройка OSPF на маршрутизаторе сети 42-го
квартала}
\phantomsection\label{ux43dux430ux441ux442ux440ux43eux439ux43aux430-ospf-ux43dux430-ux43cux430ux440ux448ux440ux443ux442ux438ux437ux430ux442ux43eux440ux435-ux441ux435ux442ux438-42-ux433ux43e-ux43aux432ux430ux440ux442ux430ux43bux430}
\begin{itemize}
\tightlist
\item
  Настроен OSPF на маршрутизаторе сети 42-го квартала
\item
  Назначен Router ID \texttt{10.128.254.2}
\item
  Сети добавлены в \texttt{area\ 0}
\end{itemize}
\end{frame}

\begin{frame}
\begin{figure}

\centering{

\includegraphics[width=0.7\linewidth,height=\textheight,keepaspectratio]{image/3.png}

}

\caption{\label{fig-3}Настройка OSPF на маршрутизаторе сети 42-го
квартала}

\end{figure}%
\end{frame}

\begin{frame}[fragile]{Настройка OSPF на маршрутизаторе сети общежития}
\phantomsection\label{ux43dux430ux441ux442ux440ux43eux439ux43aux430-ospf-ux43dux430-ux43cux430ux440ux448ux440ux443ux442ux438ux437ux430ux442ux43eux440ux435-ux441ux435ux442ux438-ux43eux431ux449ux435ux436ux438ux442ux438ux44f}
\begin{itemize}
\tightlist
\item
  Настроен OSPF на маршрутизаторе сети общежития
\item
  Назначен Router ID \texttt{10.128.254.3}
\item
  Выполнена публикация сетей в \texttt{area\ 0}
\end{itemize}
\end{frame}

\begin{frame}
\begin{figure}

\centering{

\includegraphics[width=0.7\linewidth,height=\textheight,keepaspectratio]{image/4.png}

}

\caption{\label{fig-4}Настройка OSPF на маршрутизаторе сети общежития}

\end{figure}%
\end{frame}

\begin{frame}[fragile]{Настройка OSPF на маршрутизаторе филиала в Сочи}
\phantomsection\label{ux43dux430ux441ux442ux440ux43eux439ux43aux430-ospf-ux43dux430-ux43cux430ux440ux448ux440ux443ux442ux438ux437ux430ux442ux43eux440ux435-ux444ux438ux43bux438ux430ux43bux430-ux432-ux441ux43eux447ux438}
\begin{itemize}
\tightlist
\item
  Настроен OSPF на маршрутизаторе филиала в Сочи
\item
  Назначен Router ID \texttt{10.128.254.4}
\item
  Локальные сети добавлены в \texttt{area\ 0}
\end{itemize}
\end{frame}

\begin{frame}
\begin{figure}

\centering{

\includegraphics[width=0.7\linewidth,height=\textheight,keepaspectratio]{image/5.png}

}

\caption{\label{fig-5}Настройка OSPF на маршрутизаторе филиала в Сочи}

\end{figure}%
\end{frame}

\begin{frame}[fragile]{Создание VLAN 7 на коммутаторе провайдера}
\phantomsection\label{ux441ux43eux437ux434ux430ux43dux438ux435-vlan-7-ux43dux430-ux43aux43eux43cux43cux443ux442ux430ux442ux43eux440ux435-ux43fux440ux43eux432ux430ux439ux434ux435ux440ux430}
\begin{itemize}
\tightlist
\item
  Создан VLAN 7 \texttt{q42-sochi} на коммутаторе провайдера
\item
  Активирован интерфейс \texttt{Vlan7}
\end{itemize}
\end{frame}

\begin{frame}
\begin{figure}

\centering{

\includegraphics[width=0.7\linewidth,height=\textheight,keepaspectratio]{image/6.png}

}

\caption{\label{fig-6}Создание VLAN 7 на коммутаторе провайдера}

\end{figure}%
\end{frame}

\section{Настройка подинтерфейса VLAN 7 на маршрутизаторе 42-го
квартала}\label{ux43dux430ux441ux442ux440ux43eux439ux43aux430-ux43fux43eux434ux438ux43dux442ux435ux440ux444ux435ux439ux441ux430-vlan-7-ux43dux430-ux43cux430ux440ux448ux440ux443ux442ux438ux437ux430ux442ux43eux440ux435-42-ux433ux43e-ux43aux432ux430ux440ux442ux430ux43bux430}

\begin{frame}[fragile]{Настройка подинтерфейса VLAN 7 на маршрутизаторе
42-го квартала}
\begin{itemize}
\tightlist
\item
  Создан подинтерфейс \texttt{FastEthernet0/1.7}
\item
  Настроена инкапсуляция IEEE 802.1Q
\item
  Назначен IP-адрес \texttt{10.128.255.9/30}
\item
  Добавлено описание \texttt{sochi}
\end{itemize}
\end{frame}

\begin{frame}
\begin{figure}

\centering{

\includegraphics[width=0.7\linewidth,height=\textheight,keepaspectratio]{image/7.png}

}

\caption{\label{fig-7}Настройка подинтерфейса VLAN 7 на маршрутизаторе
42-го квартала}

\end{figure}%
\end{frame}

\begin{frame}[fragile]{Создание VLAN 7 на коммутаторе филиала в Сочи}
\phantomsection\label{ux441ux43eux437ux434ux430ux43dux438ux435-vlan-7-ux43dux430-ux43aux43eux43cux43cux443ux442ux430ux442ux43eux440ux435-ux444ux438ux43bux438ux430ux43bux430-ux432-ux441ux43eux447ux438}
\begin{itemize}
\tightlist
\item
  Создан VLAN 7 \texttt{q42-sochi} на коммутаторе филиала в Сочи
\item
  Активирован интерфейс \texttt{Vlan7}
\item
  Интерфейс переведён в состояние \texttt{up/up}
\end{itemize}
\end{frame}

\begin{frame}
\begin{figure}

\centering{

\includegraphics[width=0.7\linewidth,height=\textheight,keepaspectratio]{image/8.png}

}

\caption{\label{fig-8}Создание VLAN 7 на коммутаторе филиала в Сочи}

\end{figure}%
\end{frame}

\begin{frame}[fragile]{Настройка OSPF и подинтерфейса VLAN 7 в филиале
Сочи}
\phantomsection\label{ux43dux430ux441ux442ux440ux43eux439ux43aux430-ospf-ux438-ux43fux43eux434ux438ux43dux442ux435ux440ux444ux435ux439ux441ux430-vlan-7-ux432-ux444ux438ux43bux438ux430ux43bux435-ux441ux43eux447ux438}
\begin{itemize}
\tightlist
\item
  Настроен OSPF на маршрутизаторе филиала в Сочи
\item
  Создан подинтерфейс \texttt{FastEthernet0/0.7}
\item
  Настроена инкапсуляция IEEE 802.1Q
\item
  Назначен IP-адрес \texttt{10.128.255.10/30}
\item
  Добавлено описание \texttt{q42}
\end{itemize}
\end{frame}

\begin{frame}
\begin{figure}

\centering{

\includegraphics[width=0.7\linewidth,height=\textheight,keepaspectratio]{image/9.png}

}

\caption{\label{fig-9}Настройка OSPF и подинтерфейса VLAN 7 на
маршрутизаторе филиала в Сочи}

\end{figure}%
\end{frame}

\section{Выводы}\label{ux432ux44bux432ux43eux434ux44b}

\begin{frame}{Выводы}
\begin{itemize}
\tightlist
\item
  Выполнена настройка динамической маршрутизации OSPF
\item
  Настроен VLAN 7 сети провайдера
\item
  Созданы подинтерфейсы с IEEE 802.1Q
\item
  Подтверждена корректная работа маршрутизации между сегментами сети
\end{itemize}
\end{frame}




\end{document}
